\section{Resource Availability, FAIR Status, and Repository Contents}
\label{sec:resource-availability}

HEART is released as a public resource package rather than only as a paper
appendix. The repository contains the 103-entry benchmark corpus, the
14-rubric diagnostic codebook, 14 generic repair tools, 103
benchmark-specific repair sub-tools, five worked mini cases, the literature
index used for corpus construction, and anonymized expert A/B validation
results. The benchmark corpus is available in multiple formats:
CSV and XLSX tables in \texttt{codebook/}, a human-readable catalogue in
\texttt{datasets/}, JSON files powering the interactive website under
\texttt{docs/interactive/data/}, and the live static site at
\url{https://qingjingchen.github.io/Heart/interactive/}. The 14 rubrics have a
machine-readable version in \texttt{codebook/rubrics.csv} and
\texttt{docs/interactive/data/rubrics\_full.json}; this means the rubric
identifiers, layers, questions, anchors, calibration notes, and score
distributions are exposed as structured fields rather than only as prose.

The 103 specific repair sub-tools are listed in
\texttt{codebook/sub\_tools.csv}. This file contains primary and secondary
tool mappings for each audited benchmark, so it has 206 rows while preserving
the 103 benchmark-specific primary repair tools. The five mini cases contain
both original and HEART-revised items in
\texttt{docs/interactive/data/mini\_cases.json}. The blinded expert-validation
key and anonymized expert votes are stored under \texttt{survey/results/}; raw
Google Form exports are not committed because they contain respondent initials.

We include a resource card, datasheet, and metadata schema to support FAIR
review. \texttt{RESOURCE\_CARD.md} summarizes the resource components,
formats, DOI status, version, and reproduction instructions.
\texttt{DATASHEET.md} documents motivation, composition, inclusion criteria,
recommended uses, non-recommended uses, and maintenance. The metadata and
provenance schema in \texttt{metadata/heart\_metadata\_schema.json} defines
stable fields for benchmark entries, rubric records, repair tools, mini cases,
and expert votes. Provenance fields record source URLs, BibTeX keys, evidence
notes, verification dates, annotator roles, and review status.

The current repository version is \texttt{v0.1.0}. A DOI has not yet been
minted; the repository includes \texttt{.zenodo.json} so that an indexed
Zenodo DOI can be created for the archival release and then back-filled into
\texttt{CITATION.cff}. Until that archival deposit is complete, the GitHub
repository and GitHub Pages site serve as the public review-access locations.
The resource is reusable under CC-BY-4.0 for data, codebook, documentation,
and rubric anchors, while scripts and website code are MIT licensed. Upstream
benchmark datasets remain governed by their original licenses.

The machine-readable codebook can support LLM-assisted evidence retrieval or
draft scoring, but the HEART score tables should not be read as an automatic
benchmark-of-benchmarks. Rubric calibration was performed through human expert
reading of benchmark papers, repositories, related surveys, subsequent
related-work discussions, and author-stated limitations. This design keeps
LLM assistance available for search and drafting while preserving expert
judgment for ethically sensitive scoring and claim limitation.
